\subsection{Bornologies}

DEFINITION 1. --- A bornology on a set E is a subset B of the set of all subsets of E satisfying the following conditions (cf.~GT, X, § 1.2, Remark 2).

(B1) Every subset of a set of \(\mathfrak{B}\) belongs to \(\mathfrak{B}\) .

(B2) Every finite union of a set of \(\mathfrak{B}\) belongs to \(\mathfrak{B}\) .

We say that \(\mathfrak{B}\) is covering if every element of \(E\) is contained in a set which belongs to \(\mathfrak{B}\) , or, which is the same, if \(\mathfrak{B}\) is a cover of \(E\) .

Example. --- Let E be a metric space; the set of all bounded subsets of E (GT, IX, § 2, No.~2) is a covering bornology on E. Let G be the group of isometries of E; the set of all subsets M of G such that for every \(x \in E\) , the set M.x is a bounded subset of E, is a covering bornology on G.

If \(\mathfrak{B}\) is a bornology on a set \(E\) , a subset \(\mathfrak{B}_1\) of \(\mathfrak{B}\) is said to be a base of \(\mathfrak{B}\) if every set of \(\mathfrak{B}\) is contained in a set of \(\mathfrak{B}_1\) .

The intersection of a family of bornologies on E is a bornology; consequently for every subset S of \(\mathfrak{P}(E)\) , there exists a smallest bornology containing S; this bornology is said to be generated by S and admits as a base the set of finite unions of sets of S. If E and \(E'\) are two sets, and B (resp. \(B'\) ) a bornology on E (resp. \(E'\) ), the product bornology is the bornology on \(E \times E'\) which admits the sets \(M \times M'\) as a base, where \(M \in B\) and \(M' \in B'\) .

DEFINITION 2. --- Let E be a vector space. A bornology B on E is said to be convex, if for every \(X \in B\) and \(t \in K\) , the homothetic set tX and the convex balanced envelope \(\Gamma(X)\) (II, p.~10) of X belong to B.

If X and Y are two subsets of E, we have

\[
\mathbf {X} + \mathbf {Y} \subset 2 \Gamma (\mathbf {X} \cup \mathbf {Y})
\]

\[
\lambda \mathrm{X} \subset t \Gamma (\mathrm{X}) \quad \text { for } \quad | \lambda | \leqslant t.
\]

Consequently, if B is a convex bornology on E, if A is a bounded subset of K and if X, Y belong to B, then \(X + Y \in B\) and \(A. X \in B\) .

